\documentclass[11pt]{ocots-td}

\usepackage[lang=fr, theme=ocots, solutions=none,
            math={analysis}, institution={n7}]{ocots}

\title{TD 1 --- Suites réelles}
\shorttitle{TD 1 --- Suites réelles}
\numero{TD~1}
\date{2026--2027}
\discipline{Suites réelles (cours de démonstration)}
\promotion{ocourses}

\begin{document}

\maketitle

\section{Convergence}

\begin{exercise}[label=ex:definition]
    Montrer à l'aide de la définition que la suite $u_n = \frac{2n + 1}{n + 3}$
    converge, et préciser sa limite.
\end{exercise}

\begin{exercise}[label=ex:oscillante]
    On pose $u_n = (-1)^n$. Question préliminaire~: la suite est-elle bornée ?
    \begin{question}
        Montrer que $(u_n)$ ne converge pas.
    \end{question}
    \begin{question}
        Exhiber deux sous-suites convergentes de limites différentes.
    \end{question}
\end{exercise}

\section{Suites monotones}

\begin{exercise}[label=ex:recurrente]
    On définit $u_0 = 0$ et $u_{n+1} = \sqrt{2 + u_n}$. Indication~: raisonner
    par récurrence.
    \begin{question}
        Montrer que $0 \le u_n \le 2$ pour tout $n$.
    \end{question}
    \begin{question}
        Montrer que $(u_n)$ est croissante, puis qu'elle converge vers $2$.
    \end{question}
\end{exercise}

\end{document}
